% Pages 31–40: Einstein Model of Solids, Interacting Systems, Large Systems

% ── Continuation of the paramagnet discussion (page 31 top) ──────────────────

The multiplicity of a two-state paramagnet with $N$ total dipoles and $N_\uparrow$ aligned
dipoles is
\[
  \Omega(N_\uparrow) = \frac{N!}{N_\uparrow!\, N_\downarrow!}.
\]
Parallel dipoles carry \emph{less} (more negative) energy than antiparallel ones, because
flipping a dipole against an external field requires work.  Since the total energy is
proportional to the number of up or down dipoles, specifying the macrostate (i.e.\
$N_\uparrow$) is sufficient to determine $E_{\text{tot}}$.

% ─────────────────────────────────────────────────────────────────────────────
\section{The Einstein Model of Solids}
% ─────────────────────────────────────────────────────────────────────────────

Consider a solid modelled as a collection of quantum harmonic oscillators whose energies
come in discrete units of size $hf$.  Here $h = 6.63\times10^{-34}$~J$\cdot$s is
Planck's constant and $f = \frac{1}{2\pi}\sqrt{k_s/m} = \omega/2\pi$ is the natural
oscillation frequency.  Classically the potential energy would be the continuous function
$\frac{1}{2}k_s x^2$, but quantum mechanics forces the allowed energies to be integer
multiples of $hf$ above the zero-point energy $E_0 = \frac{1}{2}hf$.

\begin{center}
[DIAGRAM: Several identical harmonic-oscillator potential wells drawn side by side.
Horizontal lines mark the equally-spaced energy levels $E_0 = \tfrac{1}{2}hf$,
$\tfrac{3}{2}hf$, $\tfrac{5}{2}hf$, \ldots\ inside each well.  Label beneath the
group: ``Oscillators''.]
\end{center}

In a three-dimensional solid each atom sits at a lattice site and vibrates independently
along three orthogonal directions, giving $N_{\text{atoms}}$ atoms a total of
$N = 3N_{\text{atoms}}$ independent oscillators.  The central combinatorial question is:
in how many ways can $q$ indistinguishable energy units be distributed among $N$
oscillators?

\begin{formula}[Multiplicity of an Einstein Solid]
The number of microstates consistent with having exactly $q$ energy units shared among
$N$ oscillators is
\[
  \Omega(N,q) = \binom{q+N-1}{q} = \frac{(q+N-1)!}{q!\,(N-1)!}.
\]
This is the ``stars and bars'' result: we are distributing $q$ identical objects (stars)
into $N$ bins separated by $N-1$ dividers (bars).
\end{formula}

As a quick check, the arrangement $*\,|\,{*}{*}\,|\,|\,{*}{*}$ uses 3 bars and therefore
describes 4 oscillators carrying a total of 5 energy units.

% ─────────────────────────────────────────────────────────────────────────────
\section{Interacting Systems}
\label{sec:interacting}
% ─────────────────────────────────────────────────────────────────────────────

Consider two Einstein solids, A and B, placed in thermal contact.  We say they are
\emph{weakly coupled}: energy can flow between them, but the rate of inter-solid exchange
is much slower than the rate of intra-solid redistribution.  Consequently, on short
timescales the individual energies $U_A$ and $U_B$ are essentially fixed, while on long
timescales only the total $U_{\text{tot}} = U_A + U_B$ is conserved.  The macrostate of
the combined system is therefore specified by the pair $(q_A, q_B)$ with $q_A + q_B = q$
constant.

\begin{center}
[DIAGRAM: Two rectangular boxes labeled ``Solid A, $N_A$, $q_A$'' and
``Solid B, $N_B$, $q_B$'' connected by a wavy arrow labeled ``Energy''.  The pair is
enclosed in a dashed rectangle labeled ``Isolated system''.]
\end{center}

Each solid has its own multiplicity:
\[
  \Omega_A = \binom{q_A + N_A - 1}{q_A}, \qquad
  \Omega_B = \binom{q_B + N_B - 1}{q_B}.
\]
Because the two solids are statistically independent, the total multiplicity factors as
$\Omega_{\text{tot}} = \Omega_A \cdot \Omega_B$ by the generalized product rule.

\begin{definition}[Fundamental Assumption of Statistical Mechanics]
In an isolated system in thermal equilibrium, every accessible microstate is equally
probable.
\end{definition}

This assumption implies that any microscopic process $X \to Y$ occurs with the same rate
as its time-reverse $Y \to X$.  It does \emph{not}, however, mean that all
\emph{macrostates} are equally likely: a macrostate with more microstates is simply more
probable.  There is also a consistency requirement on the coupling strength---the two
solids must be coupled strongly enough to exchange energy (so that $q = q_A + q_B$ is
conserved and all macrostates are reachable), yet weakly enough that the energy-level
spacing within each oscillator is undisturbed.

\begin{center}
[DIAGRAM: Bar chart with horizontal axis $q_A$ running from $0$ to $q$ (equivalently
$q_B$ from $q$ to $0$).  The bars show $\Omega_{\text{tot}}$ peaking sharply near the
centre.  A tall central bar is labelled ``State A (probable)''; a short bar near the
right edge is labelled ``State B (improbable)''.]
\end{center}

The distribution of $\Omega_{\text{tot}}$ over macrostates peaks sharply near $q_A =
q_B = q/2$.  If the system starts in an extreme macrostate B (nearly all energy in one
solid), it is overwhelmingly likely to evolve toward the peak (state A).  We describe
this as energy (heat) flowing spontaneously from the hotter solid to the cooler one.

\begin{definition}[Second Law of Thermodynamics]
Energy in an isolated system flows spontaneously until the system reaches (or fluctuates
near) its most probable macrostate---the macrostate of maximum entropy.
\end{definition}

\subsection*{Problem 2.8}

Two Einstein solids each have $N = 10$ oscillators and share a total of $q = 20$ energy
units.

\begin{enumerate}[label=\alph*)]
  \item There are 21 macrostates, corresponding to $q_A = 0, 1, \ldots, 20$.

  \item The total number of microstates follows from summing over all macrostates, which
    by the Vandermonde identity collapses to a single binomial:
    \[
      \Omega_{\text{tot}}
      = \sum_{i=0}^{20} \binom{i+9}{i}\binom{20-i+9}{20-i}
      = \binom{39}{20} \approx 6.89\times10^{10}.
    \]

  \item All 20 units in solid A corresponds to $q_A = 20$, $q_B = 0$.  The number of
    such microstates is $\binom{29}{20}$, giving a probability
    \[
      P = \frac{\dbinom{29}{20}}{\dbinom{39}{20}}
        = \frac{1.0\times10^9}{6.89\times10^{10}}
        \approx 1.45\times10^{-2}.
    \]

  \item The most probable macrostate has $q_A = q_B = 10$, with probability
    \[
      P_{\max}
      = \frac{\dbinom{19}{10}^2}{\dbinom{39}{20}}
      = \frac{8.255\times10^9}{6.89\times10^{10}}
      \approx 0.124.
    \]

  \item Starting with all energy in solid A corresponds to starting in state B of the
    diagram above; the system will spontaneously evolve toward the uniform distribution
    (state A).
\end{enumerate}

% ─────────────────────────────────────────────────────────────────────────────
\section{Large Systems}
\label{sec:large-systems}
% ─────────────────────────────────────────────────────────────────────────────

As $N$ and $q$ grow, the multiplicity distribution $\Omega_{\text{tot}}(q_A)$ becomes
dramatically sharper relative to the total range of $q_A$.

\begin{center}
[DIAGRAM: Two bell-curve sketches side by side.  Left curve is broad and low, labelled
``$N,\,q \approx \text{few hundred}$''.  Right curve is tall and narrow, labelled
``$N,\,q \approx \text{few thousand}$''.  An arrow points from left to right.]
\end{center}

Two properties of large numbers underlie all of what follows.  First, when $\lambda \gg
a$ we have $\lambda + a \approx \lambda$ (e.g.\ $10^{23} + 23 \approx 10^{23}$).
Second, multiplying a large power by a modest factor barely changes the exponent:
$y^x \cdot y^a \approx y^x$ when $y^a$ is not astronomically large relative to $y^x$
(e.g.\ $10^{10}\cdot 10 \approx 10^{10}$,\; $2^{10}\cdot 2^3 \approx 2^{10}$).
Combined with Taylor's theorem $f(a+x)\approx f(a)+f'(a)x+\cdots$, these observations
yield two workhorses:

\begin{formula}[Logarithm Approximation]
For $|x| \ll 1$,
\[
  \ln(1+x) \approx x.
\]
This follows immediately from the Taylor series
$\ln(1+x) = x - x^2/2 + x^3/3 - \cdots$
\end{formula}

\begin{formula}[Stirling's Approximation]
For $N \gg 1$,
\[
  N! \approx N^N e^{-N} \sqrt{2\pi N},
  \qquad \text{equivalently,} \qquad
  \ln(N!) \approx N\ln N - N.
\]
\end{formula}

\subsection{Multiplicity of a Large Einstein Solid}

In the \emph{high-temperature limit} $q \gg N$ (many energy units per oscillator) we can
approximate $(q+N-1)! \approx (q+N)!$ and $(N-1)! \approx N!/(N)$, giving
\[
  \Omega(N,q) \approx \frac{(q+N)!}{q!\,N!}.
\]
Taking the logarithm and applying Stirling's approximation:
\begin{align*}
  \ln\Omega
  &= \ln\!\bigl[(q+N)!\bigr] - \ln(q!) - \ln(N!) \\
  &\approx (q+N)\ln(q+N) - (q+N) - q\ln q + q - N\ln N + N \\
  &= (q+N)\ln(q+N) - q\ln q - N\ln N.
\end{align*}
Since $q \gg N$, we expand $\ln(q+N)$ using the logarithm approximation:
\[
  \ln(q+N) = \ln q + \ln\!\left(1+\frac{N}{q}\right) \approx \ln q + \frac{N}{q}.
\]
Substituting back:
\begin{align*}
  \ln\Omega
  &= (q+N)\!\left(\ln q + \frac{N}{q}\right) - q\ln q - N\ln N \\
  &= q\ln q + N + N\ln q + \frac{N^2}{q} - q\ln q - N\ln N \\
  &= N\ln\!\left(\frac{q}{N}\right) + N + \underbrace{\frac{N^2}{q}}_{\to\,0}.
\end{align*}

\begin{formula}[Einstein Solid: High-Temperature Limit]
For $q \gg N$,
\[
  \Omega(N,q) \approx \left(\frac{eq}{N}\right)^N.
\]
\end{formula}

For arbitrary large $N$ and $q$ (not necessarily in either limit), a more careful
application of Stirling's approximation to the exact formula gives:

\begin{formula}[Einstein Solid: General Large-$N$ Approximation]
\[
  \Omega(N,q) \approx
  \frac{\left(\dfrac{q+N}{q}\right)^{\!q}\!\left(\dfrac{q+N}{N}\right)^{\!N}}
       {\sqrt{2\pi q(q+N)/N}}.
\]
\end{formula}

\subsection*{Problems}

\textbf{2.17.}
In the low-temperature limit $q \ll N$ (few energy units per oscillator), expand
$\ln(N+q)$ instead:
\[
  \ln(N+q) = \ln N + \ln\!\left(1+\frac{q}{N}\right) \approx \ln N + \frac{q}{N}.
\]
Then
\begin{align*}
  \ln\Omega
  &\approx (q+N)\!\left(\ln N + \frac{q}{N}\right) - q\ln q - N\ln N \\
  &= q\ln N + \frac{q^2}{N} + N\ln N + \frac{q^2}{N} - q\ln q - N\ln N \\
  &\approx q\ln N + q - q\ln q
   = q\ln\!\left(\frac{N}{q}\right) + q,
\end{align*}
giving
\[
  \Omega(N,q) \approx \left(\frac{eN}{q}\right)^q \qquad (q \ll N).
\]

\textbf{2.18.}
Starting from the exact formula
\[
  \Omega(N,q) = \binom{N-1+q}{q} = \frac{(N-1+q)!}{(N-1)!\,q!}
  = \frac{N}{N+q}\cdot\frac{(N+q)!}{N!\,q!},
\]
apply Stirling's approximation to each factorial:
\begin{align*}
  \Omega(N,q)
  &= \frac{N}{N+q}\cdot
     \frac{(N+q)^{N+q}\,e^{-(N+q)}\sqrt{2\pi(N+q)}}
          {N^N e^{-N}\sqrt{2\pi N}\cdot q^q e^{-q}\sqrt{2\pi q}} \\
  &= \frac{N}{N+q}\left(\frac{N+q}{N}\right)^N\!\left(\frac{N+q}{q}\right)^q
     \sqrt{\frac{N+q}{2\pi Nq}}.
\end{align*}
The prefactor $N/(N+q)$ combines with the square root to give
\[
  \Omega(N,q)
  = \frac{\left(\dfrac{N+q}{N}\right)^N\!\left(\dfrac{N+q}{q}\right)^q}
         {\sqrt{\dfrac{2\pi q(N+q)}{N}}},
\]
which is the general formula stated above.

\textbf{2.19.}
For a two-state paramagnet with $N$ total dipoles and $N_\uparrow$ aligned:
\[
  \ln\Omega
  = \ln(N!) - \ln(N_\uparrow!) - \ln\!\bigl((N-N_\uparrow)!\bigr).
\]
Apply Stirling's approximation and expand $(N-N_\uparrow)\ln(N-N_\uparrow)$ for
$N_\uparrow \ll N$:
\[
  (N-N_\uparrow)\ln(N-N_\uparrow)
  = (N-N_\uparrow)\!\left(\ln N + \ln\!\left(1-\frac{N_\uparrow}{N}\right)\right)
  \approx (N-N_\uparrow)\!\left(\ln N - \frac{N_\uparrow}{N}\right).
\]
Collecting terms:
\[
  \ln\Omega \approx N_\uparrow\ln\!\left(\frac{N}{N_\uparrow}\right) + N_\uparrow,
\]
so
\[
  \Omega(N, N_\uparrow) \approx \left(\frac{Ne}{N_\uparrow}\right)^{N_\uparrow}.
\]

% ─────────────────────────────────────────────────────────────────────────────
\section{Sharpness of the Multiplicity Function}
\label{sec:sharpness}
% ─────────────────────────────────────────────────────────────────────────────

Consider two large, interacting Einstein solids each with $N$ oscillators, sharing $q =
q_A + q_B$ total energy units in the high-temperature limit $q \gg N$.  Using the
high-temperature formula for each solid,
\[
  \Omega_{\text{tot}}
  = \left(\frac{eq_A}{N}\right)^N\!\left(\frac{eq_B}{N}\right)^N
  = \left(\frac{e}{N}\right)^{2N}(q_A q_B)^N.
\]
This is maximised when $q_A q_B$ is maximised subject to $q_A + q_B = q$, i.e.\ at
$q_A = q_B = q/2$:
\[
  \Omega_{\max} = \left(\frac{e}{N}\right)^{2N}\!\left(\frac{q}{2}\right)^{2N}
               = \left(\frac{eq}{2N}\right)^{2N}.
\]

To quantify how sharply $\Omega_{\text{tot}}$ peaks, write $q_A = q/2 + x$ and
$q_B = q/2 - x$ for small $x \ll q$, so that $q_A q_B = (q/2)^2 - x^2$.  Then
\begin{align*}
  \ln\!\left[(q_A q_B)^N\right]
  &= N\ln\!\left[\left(\frac{q}{2}\right)^2\!\left(1 - \left(\frac{2x}{q}\right)^2\right)\right] \\
  &\approx N\!\left[\ln\!\left(\frac{q}{2}\right)^2 - \left(\frac{2x}{q}\right)^2\right],
\end{align*}
using $\ln(1-\epsilon)\approx -\epsilon$ for small $\epsilon = (2x/q)^2$.  Therefore

\begin{formula}[Gaussian Form of the Multiplicity Peak]
Near the maximum,
\[
  \Omega_{\text{tot}}(x) = \Omega_{\max}\, \exp\!\left[-N\!\left(\frac{2x}{q}\right)^2\right].
\]
The distribution is a Gaussian in the displacement $x = q_A - q/2$.
\end{formula}

\subsection*{Gaussian Width and the Thermodynamic Limit}

The $1/e$ half-width of the Gaussian is found by setting the exponent equal to $-1$:
\[
  N\!\left(\frac{2x}{q}\right)^2 = 1
  \;\Longrightarrow\; x = \pm\frac{q}{2\sqrt{N}}
  \;\Longrightarrow\; \Delta x = \frac{q}{\sqrt{N}}.
\]
The \emph{relative} width of the peak compared to the full range $q$ is
\[
  \frac{\Delta x}{q} = \frac{1}{\sqrt{N}}.
\]
For a macroscopic solid $N \sim 10^{23}$, so $\Delta x / q \sim 10^{-11.5}$---the peak
is vanishingly narrow.  In the \emph{thermodynamic limit} ($N \to \infty$), random
fluctuations away from the most probable macrostate are so improbable that they are
unmeasurable in practice.

\subsection*{Problem 2.22}

Two Einstein solids each have $N$ oscillators and $q = N$ energy units each (so $2N$
total energy units shared between $2N$ oscillators).

\begin{enumerate}[label=\alph*)]
  \item There are $2N+1$ macrostates (one for each value of $q_A$ from $0$ to $2N$).

  \item The total multiplicity, using the general formula with $2N$ oscillators and $2N$
    energy units, is
    \[
      \Omega(2N, 2N)
      = \frac{\left(\dfrac{4N}{2N}\right)^{2N}\!\left(\dfrac{4N}{2N}\right)^{2N}}
             {\sqrt{2\pi\cdot\dfrac{(2N)(4N)}{2N}}}
      = \frac{2^{4N}}{\sqrt{8\pi N}}.
    \]

  \item The maximum multiplicity occurs at $q_A = q_B = N$ (each solid has $N$
    oscillators and $N$ energy units), giving
    \[
      \Omega_{\max} = \bigl[\Omega(N,N)\bigr]^2 = \frac{2^{4N}}{4\pi N}.
    \]

  \item The fraction of all microstates that lie at the peak is
    \[
      \frac{\Omega_{\max}}{\Omega_{\text{tot}}}
      = \frac{2^{4N}/(4\pi N)}{2^{4N}/\sqrt{8\pi N}}
      = \frac{\sqrt{8\pi N}}{4\pi N}
      = \frac{1}{\sqrt{2\pi N}}.
    \]
    The Gaussian width in units of $q_A$ is $\Delta x = q/\sqrt{N} = 2N/\sqrt{N} =
    2\sqrt{N}$, so the fraction of the $q_A$ axis covered by the peak is $2\sqrt{N}/
    (2N) = 1/\sqrt{N}$.  The probability contained in the peak is roughly
    \[
      \frac{\Omega_{\max}}{\Omega_{\text{tot}}} \times \Delta x
      \;\sim\; \frac{1}{\sqrt{2\pi N}} \times \sqrt{N}
      \sim \sqrt{\frac{1}{2\pi}},
    \]
    consistent with the Gaussian normalisation.

    \begin{center}
    [DIAGRAM: Sharp symmetric peak of $\Omega_{\text{tot}}$ vs.\ $q_A$ over the range
    $[0,2N]$.  The peak width $\Delta x$ is labelled and is visually much narrower than
    the total range.  Peak height labelled $\Omega_{\max}$.]
    \end{center}
\end{enumerate}

\subsection*{Problem 2.23}

\begin{enumerate}[label=\alph*)]
  \item The total number of microstates of a two-state paramagnet with $N$ dipoles
    (all macrostates equally weighted) is
    \[
      \Omega_{\text{tot}}
      = \binom{N}{N/2}
      = \frac{N!}{\left[\left(\tfrac{N}{2}\right)!\right]^2}.
    \]
    Applying Stirling's approximation:
    \[
      \Omega_{\text{tot}}
      = \frac{N^N e^{-N}\sqrt{2\pi N}}
             {\left[\left(\tfrac{N}{2}\right)^{N/2} e^{-N/2}\sqrt{\pi N}\right]^2}
      = \frac{N^N e^{-N}\sqrt{2\pi N}}
             {(N/2)^N e^{-N}\,\pi N}
      = 2^N\sqrt{\frac{2}{\pi N}}.
    \]
\end{enumerate}
